%!TEX root = ../CC6.tex
\subsection{First order separable ODEs}
A \textbf{differential equation} is an equation that relates an unknown function (the dependent variable) to one or more of its derivatives. The \textbf{order} of a differential equation is the order of the highest derivative that appears in the equation. A function that satisfies the differential equation is called a \textbf{solution} and the process of finding a solution is called \textbf{solving} the differential equation. If the dependent variable depends on only a single unknown variable (the independent variable) then the differential equation is called an \textbf{ordinary differential equation} (ODE). We shall deal only with ODEs. 

Differential equations for functions with several independent variables and involving partial derivatives are called \textbf{partial differential equations} (PDEs). These are generally much harder to solve and wont be discussed in this course.\\

We shall usually discuss ODEs using $y$ to denote the dependent variable and $x$ for the independent variable, so solving the ODE involves finding $y(x).$\\

The generic first order ODE may be written in the form

\begin{equation}
    \label{eq:GFOO}
    \frac{dy}{dx}=f(x,y). \tag{$\star$}
\end{equation}

Generically this will have a one-parameter family of solutions ie. the \textbf{general solution} will contain an arbitrary constant. Assigning a particular value to this constant gives a \textbf{particular solution}. A particular solution will be singled out by the assignment of an \textbf{initial value} ie. requiring $y(x_0)=y_0,$ for given $x_0$ and $y_0.$ In this case the ODE is called an \textbf{initial value problem} (IVP). To solve an IVP, first solve the ODE to find the general solution and then determine the value of the constant so that the IVP is also satisfied.\\ 
 
Depending upon the form of the function $f(x,y)$ there are various methods that can be used to solve certain kinds of ODEs. We shall discuss these in turn.\\

The ODE \eqref{eq:GFOO} is \textbf{separable} if $f(x,y)=X(x)Y(y).$ 

It can be solved by direct integration as
\[\int \frac{1}{Y(y)}\,dy=\int X(x)\,dx.\]

\begin{example}
    Solve the IVP \[\frac{dy}{dx}=xe^{y-x} \quad
    \text{with} \quad y(0)=0.\]
    \begin{equation*}
        \begin{gathered}
            y'=xe^{-x}e^y, \quad \int e^{-y}\,dy=\int xe^{-x}\,dx, \quad e^{-y}=-xe^{-x}+\int e^{-x}\,dx=-xe^{-x}-e^{-x}+c \\
            y=-\log(e^{-x}(1+x)-c). \quad y(0)=0=\log(1-c), \quad \text{so}\quad c=0. \\
            y=-\log(e^{-x}(1+x))=x-\log(1+x).
        \end{gathered}
    \end{equation*}
\end{example}

\begin{example}
    Solve \[y'=\frac{2x(1+y^2)}{(1+x^2)^2}.\]
    \[\int \frac{1}{1+y^2}\,dy=\int \frac{2x}{(1+x^2)^2}\,dx, \quad \tan^{-1}y=-\frac{1}{1+x^2}-c, \quad y=-\tan\bigg(\frac{1}{1+x^2}+c\bigg).\]
\end{example}

\begin{example}
    Solve the IVP \[y'=\frac{x^2y-y}{1+y}, \quad y(3)=1.\]
    \begin{equation*}
        \begin{gathered}
            \int \frac{1+y}{y}\,dy=\int (x^2-1)\, dx, \quad y+\log|y|=\frac{x^3}{3}-x+c.\quad \text{Put}\ x=3\ \text{and} \ y=1 \\
            1=6+c, \quad c=-5, \quad \text{hence} \quad y+\log|y|=\frac{x^3}{3}-x-5.
        \end{gathered}
    \end{equation*}
    
    In this case the final solution can only be written in implicit form ie. it cannot be written in an explicit form $y=...$ where the right hand side is a function of $x$ only.
\end{example}

\newpage
